\documentclass{article}
\usepackage{lmodern}
\usepackage{amsmath}
\usepackage{listings}
\usepackage{fullpage}
\usepackage{parskip}
\usepackage{xcolor}
\lstset{
	basicstyle=\ttfamily,
	columns=fullflexible,
	frame=single,
	breaklines=true,
	postbreak=\mbox{\textcolor{red}{$\hookrightarrow$}\space},
}
\begin{document}
\title{Experiment `p\_4\_5\_analyse\_floats\_seq' Results}
\author{\today}
\date{}
\maketitle
\textbf{Experiment outcome: }SUCCESS
\\\textbf{Bad responses: }0
\\\textbf{Responses containing}~\texttt{assume}~\textbf{: }0
\\\textbf{Responses containing}~\texttt{expect}~\textbf{: }0
\\\textbf{Resolution attempts: }2
\\\textbf{Hard fails (resolution): }0
\\\textbf{Soft fails (resolution): }0
\\\textbf{Verification attempts: }2
\section*{Problem Specification}
\textbf{Problem name: }p\_4\_5\_analyse\_floats\_seq
\\\textbf{Natural language statement: }Write a method that takes a sequence of floatingpoint numbers and returns: - the average of the values. - the smallest of the values. - the largest of the values. - the range, that is the difference between the smallest and largest.
\\\textbf{Method signature: }p\_4\_5\_analyse\_floats\_seq(values: seq<real>) returns (average: real, min: real, max: real, range: real)
\subsection*{Ensures}
\begin{itemize}
\item \texttt{average == (sum(values) / |values| as real)}
\item \texttt{forall v :: v in values ==> min <= v}
\item \texttt{forall v :: v in values ==> max >= v}
\item \texttt{range == max - min}
\item \texttt{exists v :: v in values \&\& min == v}
\item \texttt{exists v :: v in values \&\& max == v}
\end{itemize}
\subsection*{Requires}
\begin{itemize}
\item \texttt{0 < |values|}
\end{itemize}
\subsection*{Functional Code Given}
\begin{lstlisting}
function sum(values: seq<real>): real
{
  if |values| == 0 then 0.0
  else values[0] + sum(values[1..])
}
\end{lstlisting}
\clearpage
\section*{GenAI interactions}
Below you will find all interactions between the `user' (program) and the `assistant' (OpenAI).
\subsection*{Program $\rightarrow$ GenAI}
\begin{lstlisting}
You are given the following task to perform in Dafny:

Write a method that takes a sequence of floatingpoint numbers and returns: - the average of the values. - the smallest of the values. - the largest of the values. - the range, that is the difference between the smallest and largest.

The signature should be:

p_4_5_analyse_floats_seq(values: seq<real>) returns (average: real, min: real, max: real, range: real)

The method should respect the following contract:

requires 0 < |values|, ensures average == (sum(values) / |values| as real), ensures forall v :: v in values ==> min <= v, ensures forall v :: v in values ==> max >= v, ensures range == max - min, ensures exists v :: v in values && min == v, ensures exists v :: v in values && max == v

The contract uses the following dafny code:

function sum(values: seq<real>): real
{
  if |values| == 0 then 0.0
  else values[0] + sum(values[1..])
}

These function(s) must not be used in your implementation of the method.

Produce and show only the Dafny body of this method, including the curly braces that surround it. Do not show the signature nor contract. You must not use 'assume' nor 'expect' anywhere in your code.
\end{lstlisting}
\subsection*{GenAI $\rightarrow$ Program}
\textbf{System fingerprint: }null
\\\textbf{ID: }chatcmpl-DmfXWRJi8rVfUEisnqGXaT038NUd4
\\\textbf{Model: }gpt-5.5-2026-04-23
\\\textbf{Created at: }1780491794
\\\textbf{Finish reason: }stop
\\\textbf{Usage: }624 tokens in, and 9146 tokens out
\begin{lstlisting}
{
  if |values| == 1 {
    average := values[0];
    min := values[0];
    max := values[0];
    range := max - min;

    assert values == [values[0]];
    assert values[0] in values;

    forall v | v in values
      ensures min <= v
    {
      assert v == values[0];
    }

    forall v | v in values
      ensures max >= v
    {
      assert v == values[0];
    }

    assert exists v :: v in values && min == v;
    assert exists v :: v in values && max == v;
  } else {
    assert 1 < |values|;
    var tail := values[1..];
    assert 0 < |tail|;
    assert |values| == |tail| + 1;
    assert values == [values[0]] + tail;
    assert forall v :: (v in values) <==> (v == values[0] || v in tail);

    var tailAverage, tailMin, tailMax, tailRange := p_4_5_analyse_floats_seq(tail);
    assert tailRange == tailMax - tailMin;

    var tailLength := |tail| as real;
    assert tailLength != 0.0;
    assert (|values| as real) == tailLength + 1.0;

    average := (values[0] + tailAverage * tailLength) / (|values| as real);

    if values[0] <= tailMin {
      min := values[0];
      assert values[0] in values;
      assert exists v :: v in values && min == v;
    } else {
      min := tailMin;
      ghost var w: real :| w in tail && tailMin == w;
      assert w in values;
      assert exists v :: v in values && min == v;
    }

    if values[0] >= tailMax {
      max := values[0];
      assert values[0] in values;
      assert exists v :: v in values && max == v;
    } else {
      max := tailMax;
      ghost var w: real :| w in tail && tailMax == w;
      assert w in values;
      assert exists v :: v in values && max == v;
    }

    forall v | v in values
      ensures min <= v
    {
      if v == values[0] {
        if values[0] <= tailMin {
          assert min == values[0];
        } else {
          assert min == tailMin;
          assert tailMin < values[0];
        }
      } else {
        assert v in tail;
        assert tailMin <= v;
        if values[0] <= tailMin {
          assert min == values[0];
          assert values[0] <= v;
        } else {
          assert min == tailMin;
        }
      }
    }

    forall v | v in values
      ensures max >= v
    {
      if v == values[0] {
        if values[0] >= tailMax {
          assert max == values[0];
        } else {
          assert max == tailMax;
          assert tailMax > values[0];
        }
      } else {
        assert v in tail;
        assert tailMax >= v;
        if values[0] >= tailMax {
          assert max == values[0];
          assert values[0] >= v;
        } else {
          assert max == tailMax;
        }
      }
    }

    range := max - min;
  }
}
\end{lstlisting}
\subsection*{Program $\rightarrow$ GenAI}
\begin{lstlisting}
When using dafny verify, the below error is emitted and verify fails:

problem_data/p_4_5_analyse_floats_seq/program.dfy(15,0): Error: a postcondition could not be proved on this return path
   |
15 | {
   | ^

problem_data/p_4_5_analyse_floats_seq/program.dfy(9,17): Related location: this is the postcondition that could not be proved
  |
9 | 	ensures average == (sum(values) / |values| as real)
  |                  ^^


Dafny program verifier finished with 2 verified, 1 error

Correct the error in code by altering only the method body. Produce and show only the Dafny body, including the curly braces that surround it. Do not show the signature nor contract. You must not use 'assume' nor 'expect' anywhere in your code.
\end{lstlisting}
\subsection*{GenAI $\rightarrow$ Program}
\textbf{System fingerprint: }null
\\\textbf{ID: }chatcmpl-DmfaOLkiRTAZvF8dgtNorfoAqYN5I
\\\textbf{Model: }gpt-5.5-2026-04-23
\\\textbf{Created at: }1780491972
\\\textbf{Finish reason: }stop
\\\textbf{Usage: }1594 tokens in, and 4558 tokens out
\begin{lstlisting}
{
  if |values| == 1 {
    average := values[0];
    min := values[0];
    max := values[0];
    range := max - min;

    assert values == [values[0]];
    assert values[1..] == [];
    assert sum(values[1..]) == 0.0;
    assert sum(values) == values[0] + sum(values[1..]);
    assert sum(values) == values[0];
    assert (|values| as real) == 1.0;
    assert average == sum(values) / (|values| as real);

    assert values[0] in values;

    forall v | v in values
      ensures min <= v
    {
      assert v == values[0];
    }

    forall v | v in values
      ensures max >= v
    {
      assert v == values[0];
    }

    assert exists v :: v in values && min == v;
    assert exists v :: v in values && max == v;
  } else {
    assert 1 < |values|;
    var tail := values[1..];
    assert 0 < |tail|;
    assert |values| == |tail| + 1;
    assert values == [values[0]] + tail;
    assert forall v :: (v in values) <==> (v == values[0] || v in tail);

    var tailAverage, tailMin, tailMax, tailRange := p_4_5_analyse_floats_seq(tail);
    assert tailAverage == sum(tail) / (|tail| as real);
    assert tailRange == tailMax - tailMin;

    var tailLength := |tail| as real;
    assert tailLength == (|tail| as real);
    assert tailLength != 0.0;
    assert (|values| as real) == tailLength + 1.0;

    average := (values[0] + tailAverage * tailLength) / (|values| as real);

    calc {
      tailAverage * tailLength;
      == {
        assert tailAverage == sum(tail) / tailLength;
      }
      (sum(tail) / tailLength) * tailLength;
      == {
        assert tailLength != 0.0;
      }
      sum(tail);
    }

    assert values[1..] == tail;
    assert sum(values) == values[0] + sum(values[1..]);
    assert sum(values) == values[0] + sum(tail);

    calc {
      average;
      == {
      }
      (values[0] + tailAverage * tailLength) / (|values| as real);
      == {
        assert tailAverage * tailLength == sum(tail);
      }
      (values[0] + sum(tail)) / (|values| as real);
      == {
        assert sum(values) == values[0] + sum(tail);
      }
      sum(values) / (|values| as real);
    }

    if values[0] <= tailMin {
      min := values[0];
      assert values[0] in values;
      assert exists v :: v in values && min == v;
    } else {
      min := tailMin;
      ghost var w: real :| w in tail && tailMin == w;
      assert w in values;
      assert exists v :: v in values && min == v;
    }

    if values[0] >= tailMax {
      max := values[0];
      assert values[0] in values;
      assert exists v :: v in values && max == v;
    } else {
      max := tailMax;
      ghost var w: real :| w in tail && tailMax == w;
      assert w in values;
      assert exists v :: v in values && max == v;
    }

    forall v | v in values
      ensures min <= v
    {
      if v == values[0] {
        if values[0] <= tailMin {
          assert min == values[0];
        } else {
          assert min == tailMin;
          assert tailMin < values[0];
        }
      } else {
        assert v in tail;
        assert tailMin <= v;
        if values[0] <= tailMin {
          assert min == values[0];
          assert values[0] <= v;
        } else {
          assert min == tailMin;
        }
      }
    }

    forall v | v in values
      ensures max >= v
    {
      if v == values[0] {
        if values[0] >= tailMax {
          assert max == values[0];
        } else {
          assert max == tailMax;
          assert tailMax > values[0];
        }
      } else {
        assert v in tail;
        assert tailMax >= v;
        if values[0] >= tailMax {
          assert max == values[0];
          assert values[0] >= v;
        } else {
          assert max == tailMax;
        }
      }
    }

    range := max - min;
  }
}
\end{lstlisting}
\section*{Final Program}
\begin{lstlisting}
function sum(values: seq<real>): real
{
  if |values| == 0 then 0.0
  else values[0] + sum(values[1..])
}

method p_4_5_analyse_floats_seq(values: seq<real>) returns (average: real, min: real, max: real, range: real)
	requires 0 < |values|
	ensures average == (sum(values) / |values| as real)
	ensures forall v :: v in values ==> min <= v
	ensures forall v :: v in values ==> max >= v
	ensures range == max - min
	ensures exists v :: v in values && min == v
	ensures exists v :: v in values && max == v
{
  if |values| == 1 {
    average := values[0];
    min := values[0];
    max := values[0];
    range := max - min;

    assert values == [values[0]];
    assert values[1..] == [];
    assert sum(values[1..]) == 0.0;
    assert sum(values) == values[0] + sum(values[1..]);
    assert sum(values) == values[0];
    assert (|values| as real) == 1.0;
    assert average == sum(values) / (|values| as real);

    assert values[0] in values;

    forall v | v in values
      ensures min <= v
    {
      assert v == values[0];
    }

    forall v | v in values
      ensures max >= v
    {
      assert v == values[0];
    }

    assert exists v :: v in values && min == v;
    assert exists v :: v in values && max == v;
  } else {
    assert 1 < |values|;
    var tail := values[1..];
    assert 0 < |tail|;
    assert |values| == |tail| + 1;
    assert values == [values[0]] + tail;
    assert forall v :: (v in values) <==> (v == values[0] || v in tail);

    var tailAverage, tailMin, tailMax, tailRange := p_4_5_analyse_floats_seq(tail);
    assert tailAverage == sum(tail) / (|tail| as real);
    assert tailRange == tailMax - tailMin;

    var tailLength := |tail| as real;
    assert tailLength == (|tail| as real);
    assert tailLength != 0.0;
    assert (|values| as real) == tailLength + 1.0;

    average := (values[0] + tailAverage * tailLength) / (|values| as real);

    calc {
      tailAverage * tailLength;
      == {
        assert tailAverage == sum(tail) / tailLength;
      }
      (sum(tail) / tailLength) * tailLength;
      == {
        assert tailLength != 0.0;
      }
      sum(tail);
    }

    assert values[1..] == tail;
    assert sum(values) == values[0] + sum(values[1..]);
    assert sum(values) == values[0] + sum(tail);

    calc {
      average;
      == {
      }
      (values[0] + tailAverage * tailLength) / (|values| as real);
      == {
        assert tailAverage * tailLength == sum(tail);
      }
      (values[0] + sum(tail)) / (|values| as real);
      == {
        assert sum(values) == values[0] + sum(tail);
      }
      sum(values) / (|values| as real);
    }

    if values[0] <= tailMin {
      min := values[0];
      assert values[0] in values;
      assert exists v :: v in values && min == v;
    } else {
      min := tailMin;
      ghost var w: real :| w in tail && tailMin == w;
      assert w in values;
      assert exists v :: v in values && min == v;
    }

    if values[0] >= tailMax {
      max := values[0];
      assert values[0] in values;
      assert exists v :: v in values && max == v;
    } else {
      max := tailMax;
      ghost var w: real :| w in tail && tailMax == w;
      assert w in values;
      assert exists v :: v in values && max == v;
    }

    forall v | v in values
      ensures min <= v
    {
      if v == values[0] {
        if values[0] <= tailMin {
          assert min == values[0];
        } else {
          assert min == tailMin;
          assert tailMin < values[0];
        }
      } else {
        assert v in tail;
        assert tailMin <= v;
        if values[0] <= tailMin {
          assert min == values[0];
          assert values[0] <= v;
        } else {
          assert min == tailMin;
        }
      }
    }

    forall v | v in values
      ensures max >= v
    {
      if v == values[0] {
        if values[0] >= tailMax {
          assert max == values[0];
        } else {
          assert max == tailMax;
          assert tailMax > values[0];
        }
      } else {
        assert v in tail;
        assert tailMax >= v;
        if values[0] >= tailMax {
          assert max == values[0];
          assert values[0] >= v;
        } else {
          assert max == tailMax;
        }
      }
    }

    range := max - min;
  }
}
\end{lstlisting}
\section*{Total Token Usage}
\textbf{Input tokens: }2218
\\\textbf{Output tokens: }13704
\\\textbf{Reasoning tokens: }11445
\\\textbf{Sum of `total tokens': }15922
\section*{Experiment Timings}
\textbf{Overall Experiment} started at 1780491793847, ended at 1780492060642, lasting 266795ms (266.80 seconds)
\\\textbf{Iteration \#1} started at 1780491793847, ended at 1780491972601, lasting 178754ms (178.75 seconds)
\\\textbf{Iteration \#2} started at 1780491972602, ended at 1780492060642, lasting 88040ms (88.04 seconds)
\end{document}
